\chapter{Passerelle d'ordres et gestion des ordres}\label{ch:ll:order-gateway-and-order-management}

Une stratégie annule un ordre d'achat au repos de 100 lots et, dans la même microseconde, envoie son remplacement un tick plus bas. En
chemin vers la bourse, l'annulation croise une exécution~: un vendeur vient de prendre l'ancien ordre. La bourse rejette l'annulation comme
trop tard, accepte le nouvel ordre, et un second vendeur le prend à son tour. Pendant les quelques microsecondes qui séparent la décision de la stratégie et le
rapport d'exécution, la firme détenait deux fois la position qu'elle croyait détenir, et une passerelle qui ne comptabilisait que les ordres qu'elle savait actifs
laisserait la stratégie en envoyer un troisième. Rien dans cette histoire n'est un bogue au sens ordinaire~: chaque message était valide,
et chaque partie a fait ce que le protocole prévoit. C'est ce qui arrive lorsque deux machines séparées par un réseau modifient chacune l'état d'un ordre
sans attendre l'autre. Ce chapitre construit le composant qui maintient une vue fidèle des ordres de la firme pendant que
cela se produit~: la \omterm{def:ll:order-gateway-and-order-management:gateway}{passerelle d'ordres}, avec un automate d'états qui dispose de noms pour «~j'ai demandé et je n'ai pas encore de réponse~», un décompte d'exposition
qui suppose le pire pour chaque ordre encore en vol, un limiteur qui préserve le budget de messages de la plateforme, et un
rapprochement avec la copie de la journée que tient la bourse.

\section{L'automate d'états de l'ordre}

\begin{definition}[Passerelle d'ordres, automate d'états de l'ordre]\label{def:ll:order-gateway-and-order-management:gateway}
Une \emph{passerelle d'ordres}\index{passerelle d'ordres} est le composant situé entre les stratégies d'une firme et la session de saisie des ordres d'une plateforme~:
elle transforme les demandes des stratégies en messages de la plateforme, les rapports de la plateforme en états d'ordres, et refuse ce que les limites de la firme
ou les règles de la plateforme interdisent. Un \emph{automate d'état de l'ordre}\index{automate d'état de l'ordre} est l'ensemble des états dans lesquels un ordre peut
se trouver, tel que la firme le voit, avec les événements qui le font passer d'un état à l'autre~; la vue de la firme comprend les états \emph{en attente},
atteints lorsque la firme a envoyé une demande et n'a pas encore reçu sa réponse.
\end{definition}

L'ordre côté plateforme a peu d'états~: actif, partiellement exécuté, puis exécuté, annulé ou expiré (simulateur du livre~10, son
\texttt{PROTOCOL.md}). Celui de la firme en a davantage, parce que la firme est toujours un peu en retard. Après l'envoi d'un nouvel ordre, celui-ci est
\emph{nouvel ordre en attente} jusqu'à l'arrivée de l'accusé de réception~; après une demande d'annulation ou de remplacement, il est \emph{annulation en attente} ou
\emph{remplacement en attente} jusqu'à l'arrivée de la réponse, et entre-temps l'ordre peut encore être exécuté. FIX désigne les mêmes états dans son
champ de statut d'ordre (\emph{PendingNew}, \emph{PendingCancel}, \emph{PendingReplace}, à côté de \emph{New}, \emph{PartiallyFilled},
\emph{Filled}, \emph{Canceled} et \emph{Rejected}~; chapitre~15)~: le protocole admet lui-même que les deux côtés divergent pendant un
certain temps. La \Cref{fig:ll:order-gateway-and-order-management:states} dessine l'automate de la firme.

\begin{omfigure}
\begin{tikzpicture}[font=\scriptsize,>=Stealth,st/.style={draw,rounded corners=3pt,minimum width=1.9cm,minimum height=0.62cm,
  align=center,fill=omDef!12},pend/.style={st,fill=omMeth!18,dashed},term/.style={st,fill=gray!12,double}]
\node[pend] (pn) at (0,0) {nouvel ordre\\ en attente};
\node[st] (live) at (3.3,0) {actif};
\node[st] (part) at (6.6,0) {partiel};
\node[term] (fil) at (9.9,0) {exécuté};
\node[term] (rej) at (0,-2.4) {rejeté};
\node[pend] (pc) at (3.3,-2.4) {annulation\\ en attente};
\node[pend] (pr) at (6.6,-2.4) {remplacement\\ en attente};
\node[term] (cxl) at (9.9,-2.4) {annulé};
\draw[->] (pn) -- node[above]{accusé} (live);
\draw[->] (live) -- node[above]{exécution} (part);
\draw[->] (part) -- node[above,align=center]{dernière\\ exécution} (fil);
\draw[->] (pn) -- node[left]{rejet} (rej);
\draw[->] (live) -- node[left,align=right]{annulation\\ envoyée} (pc);
\draw[->] (part) -- node[right,align=left]{remplacement\\ envoyé} (pr);
\draw[->] (pr.north west) -- node[below,sloped]{remplacé} (live.south east);
\draw[->] (pr) -- node[above]{annulé} (cxl);
\draw[->,omThm,thick] (pr.north east) -- node[pos=0.72,right,align=left,text=omThm]{dernière\\ exécution} (fil.south west);
\draw[->] (pc.south) -- ++(0,-0.7) -| node[pos=0.25,above]{annulé} (cxl.south);
\draw[->,draw=omThm,thick] ([xshift=-4mm]pc.south) |- node[pos=0.75,below,text=omThm]{dernière exécution, puis « trop tard »}
  (11.5,-3.9) |- (fil.east);
\draw[->,omThm,thick] (pc.west) .. controls +(-0.9,0.8) and +(-0.2,0.9) .. node[pos=0.3,above left,text=omThm,inner sep=1pt]{exécution partielle} ([xshift=-6mm]pc.north);
\end{tikzpicture}

\omcaption{La vue de la firme sur un ordre. Les cadres en pointillés sont les états en attente, atteints lorsqu'une demande est envoyée et quittés lorsque sa
réponse arrive~; en rouge, les exécutions qui atteignent un ordre après que la firme a demandé son annulation. Non représenté~: tout état ouvert peut aussi être
annulé par la plateforme de sa propre initiative (perte de connexion, suspension de cotation, expiration), et le nouvel ordre en attente peut être exécuté ou annulé avant
son accusé de réception. Seuls les états terminaux, doublés, sont définitifs~; la table du \cref{lst:ll:order-gateway-and-order-management:table} est
la machine complète.}\label{fig:ll:order-gateway-and-order-management:states}
\end{omfigure}

Le chantier écrit la machine sous forme de tableau, une ligne par état et une colonne par événement, avec une marque pour les transitions qui
ne sont pas autorisées (\cref{lst:ll:order-gateway-and-order-management:table}). Un tableau présente trois avantages sur un enchevêtrement de
conditions. Il est complet par construction~: chaque couple état-événement possède une entrée, de sorte que la question «~que se passe-t-il si une exécution
arrive pendant qu'un remplacement est en attente~?~» reçoit une réponse que quelqu'un a écrite. Il est une donnée~: la référence en Python, la passerelle C++
et la passerelle Rust en contiennent chacune une copie, et les tests vérifient que les trois copies forment le même tableau, entrée par entrée,
par rapport à \texttt{data/table.txt}. Enfin, un événement impossible est détecté et non absorbé~: un rapport que le tableau n'autorise pas
est renvoyé comme erreur, et c'est ainsi que la première version du chantier a constaté qu'il n'existait pas de flèche d'\emph{actif} vers
\emph{annulé}. La plateforme annule elle-même des ordres lorsqu'une session est perdue, lorsque la cotation est suspendue ou lorsqu'un ordre expire, et
le test de conformité contre le serveur réel du livre~10, qui coupe volontairement une connexion, a heurté l'entrée manquante dès sa première exécution.

\omcode{code/firm/ordergw/cpp/firm_ordergw.hpp}{28}{38}{L'automate d'états de l'ordre sous forme de tableau~: l'état suivant pour chaque état
et chaque événement, et une marque pour ce qui est interdit.}\label{lst:ll:order-gateway-and-order-management:table}

Un remplacement mérite un mot. Le simulateur, comme la plupart des plateformes, attribue à l'ordre remplacé un nouvel identifiant client, et la passerelle
conserve l'historique de l'ordre sous cet identifiant~; tant que le remplacement est en attente, l'ordre peut être exécuté sous son ancien identifiant, de sorte que la passerelle
conserve les deux jusqu'à l'arrivée de la réponse. Savoir si l'ordre garde sa place dans la file relève de la règle de la plateforme (le rapport indique
si la priorité a été conservée)~; la passerelle se contente de l'enregistrer.

\section{Accusés de réception, ordres en vol et courses}

\begin{definition}[Ordre en vol, course annulation--exécution]\label{def:ll:order-gateway-and-order-management:inflight}
Un \emph{ordre en vol}\index{ordre en vol} est un ordre dans un état en attente~: une demande le concernant a été envoyée et sa
réponse n'est pas arrivée, de sorte que l'ordre a pu changer chez la plateforme d'une manière que la firme ne connaît pas encore. Une \emph{course annulation--
exécution}\index{course annulation--exécution} se produit lorsqu'une exécution et une demande d'annulation ou de remplacement du même ordre se croisent~: l'exécution a lieu
chez la plateforme avant que la demande n'y arrive, et son rapport atteint la firme après que la demande a été envoyée.
\end{definition}

La \Cref{fig:ll:order-gateway-and-order-management:race} en dessine une. La stratégie envoie son annulation à l'instant $c$~; elle arrive chez
la plateforme, après une latence unidirectionnelle $L$~; chaque rapport d'exécution met $d$ pour revenir. Toute exécution dans l'intervalle de $c - d$ à $c + L$
entre en course avec l'annulation~: les exécutions antérieures à $c - d$ étaient connues au moment où la stratégie a décidé, et celles postérieures à $c + L$ ne peuvent pas se produire,
puisque l'ordre n'existe plus. Si les exécutions arrivent selon un processus de Poisson de taux $\lambda$ sur un ordre au repos, la probabilité qu'une annulation rencontre
au moins une est
\[
P(\text{course}) = 1 - e^{-\lambda (L + d)} \approx \lambda (L + d),
\]
pourvu que l'ordre soit resté au repos plus longtemps que $L + d$ avant l'annulation (un ordre annulé plus tôt a été exposé moins longtemps, et
le problème du week-end en calcule la probabilité exacte). La course n'est pas une défaillance de la plateforme ni de la passerelle. C'est la
latence du chemin, vue du point de vue de l'ordre, et les seuls moyens de la rendre plus rare sont de raccourcir le chemin ou d'annuler
moins.

\begin{omfigure}
\begin{tikzpicture}[font=\scriptsize,>=Stealth]
\draw[thick] (0,0) node[above]{stratégie et passerelle} -- (0,-4.4);
\draw[thick] (6.4,0) node[above]{plateforme} -- (6.4,-4.4);
\fill[omThm!12] (6.3,-1.1) rectangle (6.5,-2.9);
\draw[gray] (6.3,-1.1) -- (6.5,-1.1) node[pos=0,left,gray]{$c - d$};
\draw[->,omDef,thick] (0,-1.9) node[left]{annulation envoyée, $c$} -- node[above,sloped,pos=0.22]{annulation, latence $L$} (6.4,-2.9);
\draw[->,omThm,thick] (6.4,-1.6) node[right]{exécution} -- node[above,sloped,pos=0.4]{rapport d'exécution, délai $d$} (0,-2.9);
\node[left,align=right] at (0,-2.9) {exécution arrive~:\\ état annulation en attente};
\draw[->,omDef] (6.4,-2.9) -- node[below,sloped,pos=0.35]{« trop tard »} (0,-3.9);
\node[right,align=left,omDef] at (6.4,-3.25) {annulation arrive à $c + L$~:\\ rien à annuler};
\node[left] at (0,-3.9) {exécuté};
\draw[decorate,decoration={brace,amplitude=4pt}] (7.1,-1.1) -- node[right=4pt,align=left]{une exécution ici entre en course avec\\ l'annulation~: de $c - d$\\ à $c + L$} (7.1,-2.9);
\end{tikzpicture}

\omcaption{Une \omterm{def:ll:order-gateway-and-order-management:inflight}{course annulation--exécution}. L'annulation part à $c$ et arrive $L$ plus tard~; une exécution sur la plateforme entre $c - d$ et $c + L$
est rapportée après l'envoi de l'annulation. La passerelle, dans l'état \emph{annulation en attente}, applique l'exécution puis la réponse «~trop
tard~» et finit dans l'état \emph{exécuté}.}\label{fig:ll:order-gateway-and-order-management:race}
\end{omfigure}

La construction le mesure. \texttt{firm\_ordergw\_sim.py} exécute la passerelle Python face au moteur d'appariement du livre~10~: une stratégie
maintient un ordre d'achat en cours, le garde pendant une durée de moyenne \qty{1}{\milli\second} suivant une loi exponentielle, puis l'annule~; une
contrepartie lui vend au hasard, 200 fois par seconde~; chaque rapport met \qty{20}{\micro\second} pour revenir.
La \Cref{fig:ll:order-gateway-and-order-management:races} donne la fraction des annulations ayant rencontré une exécution, pour des latences unidirectionnelles
allant de \qty{10}{\micro\second} à \qty{3}{\milli\second}. À \qty{100}{\micro\second}, 2.1~\% des annulations ont connu une \omterm{def:ll:order-gateway-and-order-management:inflight}{course annulation--exécution}, contre
2.2~\% selon la formule exacte~; à \qty{10}{\micro\second}, elle vaut 0.65~\%. À \qty{1}{\milli\second}, la formule simple donnerait 18~\%
et le moteur annonce 12~\%~: les ordres sont gardés à peu près aussi longtemps que le trajet, si bien que la fenêtre est souvent plus courte que $L + d$.

\begin{omfigure}
\begin{tikzpicture}
\begin{axis}[width=11cm,height=6cm,xmode=log,xmin=8,xmax=4000,ymin=0,ymax=0.5,
  xlabel={latence unidirectionnelle de l'annulation $L$ ($\mu$s)},ylabel={fraction des annulations en course},
  tick label style={font=\small},label style={font=\small},grid=major,grid style={gray!25},
  legend pos=north west,legend style={font=\scriptsize,cells={anchor=west}},table/col sep=comma,
  yticklabel style={/pgf/number format/fixed}]
\addplot[omMeth,thick,dashed] table[x=latency_us,y=rested]{figdata/low-latency/21-order-gateway-and-order-management/race_curves.csv};
\addlegendentry{$1 - e^{-\lambda (L + d)}$ (ordres restés au moins $L + d$)}
\addplot[omDef,thick] table[x=latency_us,y=exact]{figdata/low-latency/21-order-gateway-and-order-management/race_curves.csv};
\addlegendentry{exacte, ordres maintenus \qty{1}{\milli\second} en moyenne}
\addplot[only marks,mark=*,omThm,mark size=2.2pt] table[x=latency_us,y=simulated]{figdata/low-latency/21-order-gateway-and-order-management/races.csv};
\addlegendentry{passerelle face au moteur d'appariement}
\end{axis}
\end{tikzpicture}

\omcaption{\omterm{def:ll:order-gateway-and-order-management:inflight}{Courses annulation--exécution} selon la latence de l'annulation, avec des exécutions au rythme de $\lambda = 200$ par seconde et des rapports
retardés $d = \qty{20}{\micro\second}$. Les points correspondent à la passerelle Python face au moteur d'appariement du livre~10 (deux simulations à graine fixée de
2\,000 ordres chacune par point)~; la courbe pleine est la probabilité exacte pour les ordres maintenus \qty{1}{\milli\second} en moyenne,
la courbe en tirets, la formule pour les ordres restés au moins $L + d$ au moment de leur annulation. Données~: \texttt{fig\_races.py}
(déterministe).}\label{fig:ll:order-gateway-and-order-management:races}
\end{omfigure}

Ce que fait la passerelle face à une course décide si elle est inoffensive. L'anecdote d'ouverture est la défaillance classique~: la passerelle a considéré
un ordre comme disparu lorsqu'elle a \emph{envoyé} l'annulation, si bien que le remplacement passait sous la limite de position et que les deux ordres ont été exécutés. La règle
qui l'empêche est la \emph{comptabilisation dans le pire cas}~: tant que la plateforme n'a pas répondu, un ordre compte pour tout ce qu'il pourrait encore
faire. Un ordre dont l'annulation est en attente compte toute sa quantité restante~; un ordre dont le remplacement est en attente compte la plus grande de sa quantité ancienne et de sa
nouvelle, car l'une ou l'autre peut être exécutée~; un ordre dont la création est en attente compte en totalité. L'exposition longue est alors la position plus la
quantité restante dans le pire cas de chaque ordre d'achat ouvert, et un nouvel ordre n'est envoyé que s'il reste sous la limite une fois ajouté à celle-ci. Avec une
limite de 100 lots, la passerelle de l'anecdote refuse le remplacement tant que la réponse à l'annulation n'est pas arrivée, et la position de la société ne
dépasse jamais 100. Le prix se paie en opportunités~: pendant un aller-retour, la stratégie ne peut pas utiliser la quantité qu'elle tente
de retirer.

La passerelle tient cette exposition sous forme de sommes courantes, une par côté, mises à jour à chaque changement d'ordre~: on retire de l'ordre sa
quantité restante dans le pire cas, on modifie l'ordre, puis on la réintègre (le \cref{lst:ll:order-gateway-and-order-management:change}). La vérification
avant un nouvel ordre ne coûte alors qu'une addition et une comparaison (le \cref{lst:ll:order-gateway-and-order-management:new}), quel que soit
le nombre d'ordres ouverts. Une passerelle qui recalcule l'exposition en parcourant ses ordres paie pour chacun d'eux à chaque
nouvel ordre~: la \cref{fig:ll:order-gateway-and-order-management:exposure} mesure environ \qty{0.4}{\micro\second} pour mille
ordres ouverts conservés dans un tableau et environ \qty{1.6}{\micro\second} dans une table de hachage, là où les sommes courantes coûtent moins que
ce que le chronomètre peut résoudre. Un teneur de marché qui cote quelques centaines d'instruments a des milliers d'ordres ouverts toute la journée.

\omcode{code/firm/ordergw/cpp/firm_ordergw.hpp}{218}{225}{Exposition dans le pire cas tenue en sommes courantes~: chaque modification d'un ordre
passe par \texttt{change}.}\label{lst:ll:order-gateway-and-order-management:change}

\omcode{code/firm/ordergw/cpp/firm_ordergw.hpp}{83}{98}{Un nouvel ordre~: la vérification de l'exposition dans le pire cas, le limiteur de messages,
puis le message de la plateforme écrit en place avec le codec du chapitre~16.}\label{lst:ll:order-gateway-and-order-management:new}

\begin{omfigure}
\begin{tikzpicture}
\begin{axis}[width=10cm,height=5.8cm,xmode=log,ymode=log,xmin=0.7,xmax=15000,ymin=0.5,ymax=40000,
  xlabel={ordres ouverts},ylabel={vérification de l'exposition (ns)},tick label style={font=\small},label style={font=\small},
  grid=major,grid style={gray!25},legend pos=north west,legend style={font=\scriptsize,cells={anchor=west}},
  table/col sep=comma]
\addplot[omThm,thick,mark=square*] table[x=open_orders,y=map_ns]{figdata/low-latency/21-order-gateway-and-order-management/measured_exposure.csv};
\addlegendentry{parcours d'une table de hachage}
\addplot[omMeth,thick,mark=*] table[x=open_orders,y=array_ns]{figdata/low-latency/21-order-gateway-and-order-management/measured_exposure.csv};
\addlegendentry{parcours d'un tableau}
\end{axis}
\end{tikzpicture}

\omcaption{Coût de la vérification de l'exposition dans le pire cas avant chaque nouvel ordre, recalculée en parcourant les ordres ouverts, en fonction
de leur nombre (médiane des appels chronométrés). Tenue en sommes courantes, comme dans \texttt{firm.ordergw}, elle coûte une addition, sous la
résolution du chronomètre à chaque taille. Mesurée sur un portable (Intel Core Ultra 7 155H) sous WSL2. Données~:
\texttt{bench\_gateway.py}.}\label{fig:ll:order-gateway-and-order-management:exposure}
\end{omfigure}

Le même banc d'essai mesure le travail complet de la passerelle pour chaque étape de la vie d'un ordre~: écrire un nouvel ordre, appliquer son
accusé de réception, une exécution partielle, écrire l'annulation et appliquer sa réponse, chacune décodée depuis le rapport binaire de la plateforme.
Chaque étape coûte de 4 à \qty{5}{\nano\second} à la médiane et de 5 à \qty{6}{\nano\second} au 99\up{e}
centile~: la logique de la passerelle n'est pas là où vont les microsecondes du trajet d'un ordre. Elles passent par le réseau et le noyau
(chapitre~13 et le livre~14), ainsi qu'aux contrôles de risque du chapitre suivant.

\section{Limiteurs et budgets de messages}

\begin{definition}[Seau à jetons (token bucket)]\label{def:ll:order-gateway-and-order-management:bucket}
Un \emph{seau à jetons}\index{seau à jetons} (token bucket) est un limiteur de débit qui contient jusqu'à $b$ jetons et en gagne $r$ par seconde~; un
message ne peut être envoyé que si un jeton est disponible, et son envoi en consomme un. Il autorise des rafales allant jusqu'à $b$ messages et, sur
tout intervalle de longueur $T$, au plus $b + rT$ messages.
\end{definition}

Les plateformes limitent le nombre de messages qu'une session peut envoyer~: un débit par seconde, parfois une rafale, parfois un poids par type de message
ou un rapport entre ordres et transactions (le limiteur de messages du livre~11, chapitre~10, et le simulateur du livre~10, qui rejette
l'excès avec le motif \texttt{T}). Une passerelle qui compte sur la plateforme pour faire respecter la limite n'en est informée que par les rejets,
après coup, et une plateforme peut déconnecter une session qui continue à la dépasser. La passerelle tient donc son propre budget, fixé
légèrement en dessous de celui de la plateforme, et refuse localement ce qui serait rejeté à distance. Le seau de la construction compte en
nano-jetons entiers, afin que les passerelles Python, C++ et Rust concordent jusqu'au dernier message sur le même journal
(le \cref{lst:ll:order-gateway-and-order-management:bucket}).

\omcode{code/firm/ordergw/cpp/firm_ordergw.hpp}{43}{50}{Le seau à jetons en nano-jetons entiers~: recharge pour le temps écoulé,
puis consommer un jeton ou refuser.}\label{lst:ll:order-gateway-and-order-management:bucket}

Que faire d'un message refusé relève de la stratégie, non de la passerelle. Une annulation refusée par le limiteur est le pire
des cas~: l'ordre reste sur le marché alors que la stratégie voulait le sortir. Les passerelles réservent donc une partie du budget aux annulations,
ou laissent passer les annulations avant les nouveaux ordres lorsque le budget est insuffisant~; la passerelle de la construction signale le motif de chaque refus
(limiteur, exposition, état) afin que la stratégie distingue un budget épuisé d'une limite. Le test de remplacement de la construction exécute une
stratégie qui remplace son ordre une fois par milliseconde en moyenne, face à un budget de 300 messages par seconde avec une rafale de deux~:
la passerelle refuse 111 de ses requêtes faute de budget, et la stratégie réessaie.

\section{Sessions~: connexion, numéros de séquence, annulation à la déconnexion}

La session de saisie des ordres de la plateforme (livre~10, chapitre~26) transporte les messages de la passerelle. Le simulateur du livre~10 parle SoupBinTCP~:
le client se connecte avec un nom d'utilisateur et un mot de passe~; ses messages ne sont pas numérotés, alors que les rapports du serveur le sont, numérotés
à partir de un pour chaque session, mais les numéros ne circulent pas sur le fil~: chaque côté les compte. Un client qui perd sa connexion se reconnecte
en demandant le prochain rapport qu'il n'a pas vu, et le serveur rejoue tous les rapports à partir de ce numéro. La partie du protocole qui concerne le client
est courte (\cref{lst:ll:order-gateway-and-order-management:session})~: elle compte les messages numérotés, demande le suivant
au moment de la connexion, vérifie que le serveur repart de ce numéro, envoie un \omterm{def:ll:protocols-i-fix:session}{heartbeat} après une seconde de silence, et considère le serveur comme mort
au bout de quinze secondes.

\omcode{code/firm/ordergw/firm_ordergw.py}{291}{305}{Le côté client de la session de saisie des ordres~: une connexion acceptée au
mauvais numéro de séquence est une erreur, non un détail.}\label{lst:ll:order-gateway-and-order-management:session}

Une connexion perdue est le cas où la comptabilité du pire cas prouve son utilité. Tant que la passerelle est déconnectée, elle ignore ce que font ses
ordres~: ils peuvent être en cours d'exécution, ou la plateforme peut les avoir annulés. La plupart des plateformes proposent, et beaucoup exigent, une annulation à la
déconnexion~: la plateforme annule les ordres au repos de la session lorsqu'elle détecte la perte de la connexion (le simulateur le fait
par défaut, avec le motif d'annulation \texttt{D}). La passerelle les compte encore en totalité tant que les rapports ne disent pas le contraire, car
l'annulation à la déconnexion est une fonctionnalité de la plateforme, non une garantie pour la société~: elle n'agit que sur les ordres qu'elle couvre, au moment où elle
détecte la perte, et une exécution peut la précéder. Le test de conformité fait exactement cela face au serveur en direct du livre~10 sur l'interface
loopback~: trois ordres partent, un est exécuté, la connexion est coupée, l'exposition reste à 300 lots tant que la passerelle est
absente, et après la nouvelle connexion les rapports rejoués annulent les deux autres et la ramènent à 100.

\begin{dated}{2026-09}{Annulation à la déconnexion et drop copy sur une bourse}\label{dat:ll:order-gateway-and-order-management:cme}
La documentation client de CME Group (consultée en septembre 2026) décrit le Cancel On Disconnect pour ses sessions de saisie des ordres iLink~: lorsqu'une
perte de connexion est détectée, elle annule tous les ordres au repos sur contrats à terme et sur options de la session déconnectée,
à l'exception des ordres good-till-cancelled et good-till-date, et les nouvelles sessions sont créées avec cette option activée par défaut. Son service Drop Copy
envoie en temps réel des copies des rapports d'exécution et des accusés de réception transmis sur les sessions iLink, par une voie distincte et
dédiée.
\end{dated}

\section{Drop copies et réconciliation}

\begin{definition}[Drop copy]\label{def:ll:order-gateway-and-order-management:dropcopy}
Un \emph{drop copy}\index{drop copy} est une session distincte sur laquelle la plateforme envoie des copies des rapports des sessions de négociation de la société
(au moins ses exécutions et ses annulations) au moment même où elle envoie les originaux~; elle est en lecture seule et sert de relevé indépendant
des ordres et des exécutions de la société.
\end{definition}

La vision qu'a la passerelle des ordres de la société se construit à partir des rapports reçus, appliqués dans l'ordre~; tout rapport perdu, dupliqué ou mal appliqué
rend la vision fausse, et une vision fausse ne se détecte qu'en la comparant à une autre. Le \omterm{def:ll:order-gateway-and-order-management:dropcopy}{drop copy} est cette autre
vision, venant de la plateforme, sur une autre connexion, souvent lue par un autre processus (le simulateur du livre~10 offre une session de
ce type par société). La réconciliation compare les deux~: pour chaque ordre, la quantité exécutée et s'il a été annulé
(\cref{lst:ll:order-gateway-and-order-management:reconcile}). Chaque écart est une \emph{rupture}, et une rupture s'instruit,
on ne la corrige pas en écrasant un des deux côtés~: un ordre que le \omterm{def:ll:order-gateway-and-order-management:dropcopy}{drop copy} montre exécuté et que la passerelle montre actif est un rapport perdu et une
position inconnue~; le cas inverse est un rapport appliqué deux fois.

\omcode{code/firm/ordergw/firm_ordergw.py}{189}{203}{Réconciliation avec le drop copy~: quantités exécutées et annulations,
ordre par ordre~; tout écart est une rupture.}\label{lst:ll:order-gateway-and-order-management:reconcile}

Dans la construction, chaque simulation et le test en direct se terminent sans aucune rupture, et les tests de propriétés garantissent que ce n'est pas
un coup de chance~: sur des entrelacements aléatoires d'exécutions, d'annulations, de remplacements et de latences allant de \qty{5}{\micro\second} à \qty{2}{\milli\second},
aucun ordre n'a jamais une quantité exécutée supérieure à la quantité passée, aucune quantité restante n'est négative, la position égale la somme des exécutions, et
l'exposition dans le pire cas n'est jamais inférieure à la position~; et avec une limite de 100 lots, la position face au moteur d'appariement
ne dépasse jamais 100. La société réconcilie à plusieurs rythmes~: en continu avec le \omterm{def:ll:order-gateway-and-order-management:dropcopy}{drop copy}, chaque matin avec
les registres du courtier (le livre~1, chapitre~7), et le système de gestion des ordres au-dessus de la passerelle (livre~10,
chapitre~20) avec les comptes du portefeuille.

\section{Tutoriel~: courses, limiteurs et une réconciliation}\label{tut:ll:order-gateway-and-order-management:tut}

\begin{tutorial}
\textbf{Objectif.} Faire tourner la passerelle face au moteur d'appariement, compter ses courses, rejouer son journal dans trois langages, et
la réconcilier avec le \omterm{def:ll:order-gateway-and-order-management:dropcopy}{drop copy} d'une session en direct. \textbf{État final~:}
\cref{fig:ll:order-gateway-and-order-management:races}, les rejeux identiques à \texttt{data/expected.txt}, et une réconciliation
sans aucune rupture.
\begin{tutsteps}
\item \textbf{Courses.} \texttt{python fig\_races.py} lance \texttt{firm\_ordergw\_sim.run} pour six latences et écrit les
fractions qui figurent dans la \cref{fig:ll:order-gateway-and-order-management:races}, à côté de la formule exacte de \texttt{ll\_gateway.py}.
\item \textbf{Journaux.} \texttt{make\_ordergw\_fixtures.py} enregistre deux exécutions (annulation et nouvel ordre~; annulation-remplacement sous un
budget serré) sous forme de journaux de requêtes et de rapports~; les passerelles C++ et Rust les rejouent, en refaisant chaque requête avec le même
résultat, et affichent les synthèses qu'affiche la référence Python. Le test C++ fournit aussi les rapports sous forme de messages binaires encodés,
et ne compte aucune allocation.
\item \textbf{Session.}
\texttt{tests/test\_live\_session.py} démarre le serveur en direct du livre~10, connecte une session de négociation, un \omterm{def:ll:order-gateway-and-order-management:dropcopy}{drop copy} et une contrepartie, coupe la connexion de la session de négociation, se reconnecte, et réconcilie.
\item \textbf{Coûts.} \texttt{python bench\_gateway.py} mesure chaque étape de la vie d'un ordre et le contrôle d'exposition face au
nombre d'ordres ouverts.
\end{tutsteps}
\textbf{Pour aller plus loin.} Abaisser le budget du test de remplacement à 100 messages par seconde et observer les refus~; faire
arriver les exécutions de la contrepartie par rafales plutôt qu'à un rythme constant, et comparer le nombre de courses à la formule.
\end{tutorial}

\section{Construction~: la passerelle d'ordres}\label{bld:ll:order-gateway-and-order-management:ordergw}

\begin{build}
\bfield{Rôle} Le point d'accès unique de la société à la saisie des ordres de la plateforme~: chaque ordre venant du moteur de stratégie (chapitre~20) passe
par elle, après la barrière de risque (chapitre~22)~; son journal est consigné (chapitre~23), et son état survit à une bascule (failover)
(chapitre~24).
\bfield{Interface} Référence Python \texttt{firm\_ordergw}~: \texttt{TABLE}, \texttt{transition}, \texttt{TokenBucket},
\texttt{Gateway(rate, burst, max\_long, max\_short)} avec \texttt{new}, \texttt{cancel}, \texttt{replace}, \texttt{on\_report},
\texttt{worst\_long}, \texttt{worst\_short}, \texttt{races}~; \texttt{Session}, \texttt{frames}, \texttt{reconcile},
\texttt{replay}, \texttt{summary}~; \texttt{firm\_ordergw\_sim.run} face au moteur du livre~10. C++20 \texttt{firm::ogw}~:
\texttt{Gateway(Limits)} avec \texttt{new\_order}, \texttt{cancel}, \texttt{replace} qui écrivent sur place les messages de la plateforme,
\texttt{on\_report} et \texttt{on\_wire}, \texttt{summary}. Rust \texttt{firm\_ordergw}~: \texttt{Gateway}, \texttt{replay}.
\bfield{Règles} Chaque ordre ouvert compte pour le pire qu'il peut encore faire~; rien n'est envoyé au-delà de la limite d'exposition ni du budget~;
un rapport que la table n'autorise pas est une erreur~; un ordre est conservé jusqu'à un rapport terminal~; les trois passerelles tiennent la même
table~; aucune allocation par message.
\bfield{Tests d'acceptation} \texttt{code/firm/ordergw/}~: la table identique dans trois langages~; les deux journaux rejoués jusqu'aux
synthèses versionnées en Python, C++ (à partir des champs et à partir du fil) et Rust~; tests de propriétés sur des exécutions aléatoires face au
moteur d'appariement~; la limite d'exposition tenue face au moteur~; la double position de la passerelle naïve reproduite et refusée par
la comptabilité du pire cas~; la fraction de courses à moins d'un point de la formule~; la session en direct avec annulation à la déconnexion, le rejeu
et une réconciliation sans aucune rupture.
\bfield{Pour aller plus loin} Annulation de masse (mass cancel) et cotation de masse (mass quote)~; un budget réservé aux annulations~; la couche de session en C++, adossée au
\texttt{exchsim\_client.hpp} du livre~10~; \omterm{def:ll:protocols-i-fix:session}{sessions FIX} par le moteur du chapitre~15.
\end{build}

\begin{omsources}
\item CME Group Client Systems Wiki, \emph{Cancel on Disconnect}, et \emph{Drop Copy 4.0 Service for iLink}.
\item FIX Trading Community, FIX Latest, jeu de codes \emph{OrdStatus}.
\end{omsources}

\section{Exercices}

\begin{exercise}[$\star$]\label{exo:ll:order-gateway-and-order-management:1}
Un ordre est en attente d'annulation. Une exécution portant sur toute sa quantité restante arrive, puis le rejet de l'annulation
par la plateforme, arrivé trop tard. Donner l'état de l'ordre après chaque rapport, et dire ce que la passerelle doit compter comme quantité restante entre les deux.
\end{exercise}

\begin{exercise}[$\star$]\label{exo:ll:order-gateway-and-order-management:2}
Un \omterm{def:ll:order-gateway-and-order-management:bucket}{seau à jetons} (token bucket) gagne 300 jetons par seconde et en contient au plus 2~; il est plein à l'instant 0. Des messages sont envoyés aux instants 0, 0, 0, 1, 4 et
5 millisecondes. Lesquels passent~?
\end{exercise}

\begin{exercise}[$\star$]\label{exo:ll:order-gateway-and-order-management:3}
La limite longue est de 500 lots. La position est longue de 200 lots, et deux ordres d'achat de 100 lots sont en cours. La stratégie
demande de remplacer l'un d'eux par un achat de 300. La passerelle envoie-t-elle le remplacement~? Quelle est la plus grande nouvelle quantité qu'elle enverrait~?
\end{exercise}

\begin{exercise}[$\star\star$]\label{exo:ll:order-gateway-and-order-management:4}
Des exécutions arrivent à raison de 50 par seconde sur un ordre au repos, l'annulation met \qty{300}{\micro\second} pour atteindre la plateforme et les rapports
mettent \qty{50}{\micro\second} pour revenir. Quelle fraction des annulations entre en course avec une exécution~? Combien de courses une stratégie qui annule
20\,000 ordres par jour doit-elle en attendre~?
\end{exercise}

\begin{exercise}[$\star\star$]\label{exo:ll:order-gateway-and-order-management:5}
Le \omterm{def:ll:order-gateway-and-order-management:dropcopy}{drop copy} montre deux exécutions de 60 lots pour l'ordre 17, toutes deux avec le numéro d'appariement 88\,412~; la passerelle
indique l'ordre 17 exécuté pour 60. Quel côté a tort, et que la réconciliation aurait-elle dû faire~?
\end{exercise}

\begin{exercise}[$\star\star$]\label{exo:ll:order-gateway-and-order-management:6}
Une passerelle recalcule l'exposition en parcourant une table de hachage de ses ordres ouverts avant chaque nouvel ordre. Avec mille ordres ouverts,
que cela ajoute-t-il à chaque ordre selon la \cref{fig:ll:order-gateway-and-order-management:exposure}, et comment cela se compare-t-il
au travail propre de la passerelle par message~?
\end{exercise}

\begin{exercise}[$\star\star\star$]\label{exo:ll:order-gateway-and-order-management:7}
\emph{Programmation.} Ajouter l'annulation de masse (mass cancel) à la passerelle~: une seule requête qui annule chaque ordre ouvert d'un côté. Par quels états passe chaque ordre,
que compte l'exposition jusqu'à l'arrivée des réponses, et comment le tester contre le moteur du livre~10~?
\end{exercise}

\begin{exercise}[$\star\star\star$]\label{exo:ll:order-gateway-and-order-management:8}
\emph{Trouver la faille.} «~Pour garder la table des ordres petite, notre passerelle libère l'emplacement d'un ordre dès qu'elle
envoie l'annulation~: la stratégie ne veut plus de cet ordre.~»
\end{exercise}

\section{Problème~: L'annulation qui a croisé une exécution}

\begin{problem}[{Problème du week-end --- la fenêtre dans laquelle une annulation peut se perdre}]\label{pb:ll:order-gateway-and-order-management:1}
Un teneur de marché cote des ordres d'achat de 100 lots, garde chacun en moyenne une milliseconde, puis l'annule. Les exécutions arrivent
à $\lambda = 200$ par seconde sur un ordre au repos, les rapports mettent $d = \qty{20}{\micro\second}$ pour revenir, et l'annulation met
$L$ pour atteindre la plateforme. Utiliser la \cref{fig:ll:order-gateway-and-order-management:races} (\texttt{races.csv}).

\medskip\noindent\textbf{Partie I --- La fenêtre.}
\begin{enumerate}
\item Expliquer pourquoi une exécution à la plateforme située entre $c - d$ et $c + L$ entre en course avec une annulation envoyée à $c$, et
pourquoi les exécutions hors de cet intervalle n'y entrent pas.
\item Établir $P(\text{course}) = 1 - e^{-\lambda(L + d)}$ pour un ordre resté au repos depuis au moins $L + d$ lorsque l'annulation est envoyée.
\item Calculer cette probabilité pour $L = \qty{10}{\micro\second}$, \qty{100}{\micro\second} et \qty{1}{\milli\second}.
\item Montrer qu'elle vaut à peu près $\lambda (L + d)$ lorsque ce produit est petit, et dire à partir de quelle valeur il l'est.
\end{enumerate}

\medskip\noindent\textbf{Partie II --- La mesure.}
\begin{enumerate}[resume]
\item Lire les fractions simulées aux trois mêmes latences dans les données de la figure.
\item Pourquoi la formule simple les surestime-t-elle à \qty{1}{\milli\second}~?
\item Pour un ordre tenu pendant un temps $h$, la fenêtre vaut $\min(h, L + d)$. Que devient la probabilité de course lorsque $L$ devient
très supérieur à la durée de détention~?
\item Le journal des annulations engagées (latence \qty{100}{\micro\second}) compte 6 courses sur 245 annulations. Est-ce compatible avec le
la figure~?
\end{enumerate}

\medskip\noindent\textbf{Partie III --- Les dégâts.}
\begin{enumerate}[resume]
\item Une passerelle naïve compte un ordre comme disparu dès qu'elle envoie l'annulation, et la stratégie envoie aussitôt un remplacement. Quelle
position la société peut-elle atteindre après une course, avec une limite de 100 lots~?
\item Avec $n$ courses dont tous les remplacements sont encore en cours, quelle est la pire sur-position~?
\item Que fait à la place la comptabilisation dans le pire cas, et que coûte-t-elle à la stratégie~?
\item À \qty{100}{\micro\second}, combien de courses une stratégie qui annule 40\,000 fois par jour doit-elle en attendre~?
\end{enumerate}

\medskip\noindent\textbf{Partie IV --- Le verdict.}
\begin{enumerate}[resume]
\item Énoncer le \emph{résultat nommé}~: la probabilité de course par annulation, et la pire sur-position sans
comptabilisation des \omterm{def:ll:order-gateway-and-order-management:inflight}{ordres en vol}.
\item Quelles deux grandeurs une société contrôle-t-elle pour réduire les courses, et laquelle ne contrôle-t-elle pas~?
\item Pourquoi diviser par deux la latence dans un sens ne divise-t-il pas par deux les courses lorsque $d$ est grand~?
\item Pourquoi une course n'est-elle pas une erreur à signaler à la plateforme~?
\item Pourquoi la passerelle doit-elle garder un ordre annulé dans sa table jusqu'au rapport terminal~?
\item Que change un remplacement dans l'analyse~?
\item Comment vérifier, un jour de production, que les courses observées correspondent à la formule~?
\item En une phrase~: quelle est la règle qui rend les courses inoffensives~?
\end{enumerate}
\end{problem}

\section{Questions d'entretien}

\begin{interviewq}[$\star$ \iqroles{développeur}]\label{iq:ll:order-gateway-and-order-management:1}
Par quels états passe un ordre dans une \omterm{def:ll:order-gateway-and-order-management:gateway}{passerelle d'ordres}, et pourquoi y en a-t-il davantage qu'à la bourse~?
\end{interviewq}

\begin{interviewq}[$\star\star$ \iqroles{développeur}]\label{iq:ll:order-gateway-and-order-management:2}
Vous envoyez une annulation et recevez une exécution pour le même ordre. Que s'est-il passé, et que doit faire votre système~?
\end{interviewq}

\begin{interviewq}[$\star\star$ \iqroles{développeur}]\label{iq:ll:order-gateway-and-order-management:3}
Implémentez un limiteur de débit pour une session de la bourse. Que faites-vous d'une annulation lorsque le budget est épuisé~?
\end{interviewq}

\begin{interviewq}[$\star\star$ \iqroles{développeur, trader}]\label{iq:ll:order-gateway-and-order-management:4}
Votre connexion à la bourse tombe pendant deux secondes. Que savez-vous de vos ordres, et que faites-vous
à la reconnexion~?
\end{interviewq}

\begin{interviewq}[$\star\star$ \iqroles{développeur}]\label{iq:ll:order-gateway-and-order-management:5}
Comment vérifier une limite de position en temps constant avec des milliers d'ordres ouverts~?
\end{interviewq}

\begin{interviewq}[$\star\star\star$ \iqroles{développeur, chercheur}]\label{iq:ll:order-gateway-and-order-management:6}
Concevez la couche de gestion des ordres d'un teneur de marché sur trois plateformes. Quelle est la source de vérité pour les positions, et comment
savez-vous qu'elle est juste~?
\end{interviewq}
